\section{Conclusion}
In this work, we investigated whether an expert-based correction system can
improve the cancer-related classification performance of a baseline model on the
PathMNIST dataset. The results show that this strategy can be effective when the
baseline model does not fail uniformly, but instead shows specific and repeated
confusion cases. By introducing specialized binary expert classifiers for these
confusion pairs, the final cancer expert system was able to correct many of
these errors and substantially improve the stroma classification performance.

The most important improvement is observed for stroma recall and stroma $F_2$
score. Since stroma was the weakest cancer-related class in the baseline model,
improving this class also led to stronger cancer-vs.-rest performance. This
shows that expert systems can be used to improve a general baseline classifier
in a targeted way, instead of trying to improve all classes equally. The
approach is therefore especially useful when the main objective is not only high
overall accuracy, but better performance on medically relevant classes where
false negatives are particularly problematic.

However, the presented system should be understood as a toy example rather than
a complete medical classification system. The baseline model and the expert
classifiers use relatively small CNN architectures, and only a limited amount of
hyperparameter optimization was performed. Future work could explore more
systematic hyperparameter selection for both the baseline model and the binary
experts. This could include different learning rates, regularization strategies,
augmentation settings, model sizes, and decision thresholds. A better optimized
baseline model could reduce the number of initial errors, while better optimized
experts could improve the correction step without increasing false-positive
stroma predictions.

Another limitation is the use of a fixed confidence threshold for deciding when
an expert should be applied. In this work, the threshold is manually selected
and applied equally across the relevant expert cases. Future work could
investigate more dynamic decision strategies. For example, the threshold could
be selected separately for each confusion pair, adapted based on validation-set
performance, or replaced by a learned gating mechanism that decides whether the
baseline prediction should be trusted or passed to an expert classifier.

Finally, real-world applications would likely rely on more advanced model
architectures than the small CNN used in this work. Modern vision models, such
as Vision Transformers or larger pretrained convolutional networks, could
provide stronger feature representations and better generalization.
Nevertheless, the results demonstrate the potential of combining a general
classifier with specialized expert models. Even in this simplified setting, the
expert system improved the classification of the most problematic cancer-related
class and showed that targeted correction of specific confusion cases can be a
promising strategy for improving recall-oriented medical image classification.
